% ============================================================================
% STANDALONE WRAPPER. Everything down to the marker "END OF STANDALONE
% WRAPPER" exists so that this file compiles on its own. Delete this block and
% the matching one at the end of the file to paste the body into a thesis.
% ============================================================================
% CONVENTION (default, change here): top level is \section in the article
% class, so the body drops into a thesis chapter without renaming.
\documentclass[11pt,a4paper]{article}
\usepackage[T1]{fontenc}
\usepackage[utf8]{inputenc}
\usepackage{lmodern}
\usepackage[margin=25mm]{geometry}
\usepackage{amsmath,amssymb,amsthm,booktabs,array}
\usepackage{microtype}
\usepackage[numbers,sort&compress]{natbib}
\usepackage{xurl}
\usepackage[hidelinks]{hyperref}
\newcommand{\R}{\mathbb R}
\newcommand{\col}{\operatorname{col}}
\newcommand{\rank}{\operatorname{rank}}
\newcommand{\diag}{\operatorname{diag}}
\newcommand{\refset}{\mathrm{ref}}
% CONVENTION (default, change here): theorem environments are declared in this
% deletable preamble, matching the surrounding project documents.
\newtheorem{proposition}{Proposition}
\newtheorem{lemma}{Lemma}
\theoremstyle{remark}
\newtheorem{remark}{Remark}
\title{An additional-state extension of orthogonal-by-construction\\
input-output augmentation, and its delay-state realisation}
\author{Eindhoven University of Technology}
\date{15 September 2026}
\begin{document}
\maketitle
% ================== END OF STANDALONE WRAPPER (opening) ====================

\section{Setting}\label{sec:setting}
% CONVENTION (default, change here): equations are numbered continuously and
% only where they are referred to later; all others use starred environments.
A physics-based baseline model that is linear in its physical parameters can be extended by a
learned component fitted alongside it. When the two are fitted together the learned component may
supply output behaviour that the baseline could have supplied itself, and the physical parameters
then absorb whatever remains, which is how an augmented model loses the interpretability that
motivated the baseline. The orthogonal-by-construction approach of
Gy\"or\"ok et al.~\cite{gyorok2026obc} removes the baseline-representable part of the learned
correction before that correction is added, so that on a declared reference set the two components
cannot represent the same thing.

The source construction is stated for discrete-time input-output predictors and this chapter stays
in those coordinates. The regressor is then built from past inputs and outputs, so it is available
from data without a state observer, and the finite history is itself a state, so the construction
can be examined as a state-space object without an observability hypothesis. The delay states
below are stored signals, not physical positions or velocities.

Inherited from the source are the input-output model class, the stacked basis of baseline
evaluations, the projection defining the auxiliary coefficient, and the orthogonality it enforces
on the reference set. Contributed here are the explicit delay-state realisation with its fixed
routing matrices, the extension of the model class by a learned state carrying its own update, and
the extension of the projection to the resulting enlarged state equation.

The main result is an exact reduction: isometric output insertion and zero baseline learned-state
rows make the enlarged projection equivalent to projecting the complete learned output correction,
at any rank. Additional memory enlarges the model class without requiring a separate subtraction
on its update rows. The state-level and input-output orthogonality conditions therefore coincide
in these delay coordinates. Extension to physical state-space baselines without full-state
measurement is a separate problem identified as future research
by Gy\"or\"ok et al.~\cite[Sect.~6]{gyorok2026obc}.

\section{The inherited construction}\label{sec:inherited}
The source model is a measured-output process described through a finite input-output
regressor:
\begin{equation*}
y_k=f(x_k)+e_k,\qquad
x_k=\col\bigl(y_{k-1},\ldots,y_{k-n_{\mathrm a}},
u_k,u_{k-1},\ldots,u_{k-n_{\mathrm b}}\bigr).
\end{equation*}
The regressor holds \(n_{\mathrm a}\) past outputs together with the current input and
\(n_{\mathrm b}\) past inputs. Because the realisation of Section~\ref{sec:realisation} stores the
past signals while the current input arrives from outside, we split the regressor once and keep
the two resulting arguments for the whole chapter:
\begin{equation}\label{eq:history}
s_k=\col(y_{k-1},\ldots,y_{k-n_{\mathrm a}},u_{k-1},\ldots,u_{k-n_{\mathrm b}})\in\R^{n_s},
\qquad n_s=n_{\mathrm a}n_{\mathrm y}+n_{\mathrm b}n_{\mathrm u},
\qquad \phi(s_k,u_k):=\phi(x_k).
\end{equation}
The two-argument notation retains all prediction information and uses the displayed source
ordering, with the current input before the past inputs; concatenating the history and current
input directly would put it last. Both lag orders are assumed at least one for the displayed
routing blocks; an order-one shift is zero. The baseline predictor is
linear in its parameters while its matrix function may depend nonlinearly on those signals,
\begin{equation}\label{eq:baselineIO}
\hat y_k=\phi(s_k,u_k)\theta_b,\qquad\theta_b\in\R^{n_{\theta_{\mathrm b}}},
\end{equation}
and the learned component enters additively as a correction on the output,
\begin{equation*}
\hat y_k=\phi(s_k,u_k)\theta_b+g_{y,\eta}^{(0)}(s_k,u_k).
\end{equation*}
No state or input linearisation has been performed in \eqref{eq:baselineIO}, and the learned map
is a feedforward function of the regressor with no recursion of its own.

\begin{table}[ht]
\centering\small
\begin{tabular}{ll}\toprule
Symbol & Meaning\\\midrule
\(n_{\mathrm y},\ n_{\mathrm u}\) & Output and input dimension\\
\(n_{\mathrm a},\ n_{\mathrm b}\) & Number of past output and past input samples\\
\(x_k\) & Input-output regressor, including the current input\\
\(s_k\in\R^{n_s}\) & Stored history, excluding the current input\\
\(x_{a,k}\in\R^{n_{x_a}}\) & Learned state\\
\(\xi_k\) & Combined history and learned state\\
\(\theta_b\in\R^{n_{\theta_{\mathrm b}}}\) & Baseline parameter vector\\
\(\eta\) & Parameters of the learned maps\\
\(\theta_{\mathrm{aux}}\) & Shared projection coefficient\\
\(N\) & Number of stacked evaluations\\
\(\dagger\) & Moore-Penrose pseudoinverse\\
\bottomrule\end{tabular}
\caption{Notation.}\label{tab:notation}
\end{table}

The projection is defined on \(N\) declared reference arguments, by stacking the baseline basis
and the learned correction evaluated there:
\begin{equation}\label{eq:originalstacks}
\Phi=\col_{i=1}^N\phi(s_i^{\refset},u_i^{\refset})
\in\R^{Nn_{\mathrm y}\times n_{\theta_{\mathrm b}}},
\qquad
G_\eta^{(0)}=\col_{i=1}^N g_{y,\eta}^{(0)}(s_i^{\refset},u_i^{\refset})\in\R^{Nn_{\mathrm y}}.
\end{equation}
These are function evaluations at declared arguments. Estimation histories, an auxiliary
synthetic set, or a subset of either are all admissible. Reference arguments may also be obtained
from a declared trajectory-generation procedure; exact orthogonality holds on the evaluations
used to construct the coefficient and does not automatically hold on other trajectories.
Combining distinguishability of the
baseline function with a weakly persistently exciting input sequence, the source assumes
\begin{equation}\label{eq:rank}
\rank(\Phi)=n_{\theta_{\mathrm b}},\qquad Nn_{\mathrm y}\ge n_{\theta_{\mathrm b}},
\end{equation}
which permits a normal-matrix inverse and a unique baseline least-squares coefficient. Neither
predictor evaluation nor the routing of Section~\ref{sec:realisation} requires that inverse, so
failure of \eqref{eq:rank} does not prevent realisation, although it can prevent identification of
individual baseline coefficients.

The baseline-representable part of the stacked correction is its Euclidean projection onto the
column space of the basis, and subtracting that part leaves what the baseline cannot reach.
Written at any rank with the Moore-Penrose convention, in place of the source's full-rank
inverse~\cite[Eq.~(8)]{gyorok2026obc},
\begin{equation}\label{eq:originalprojection}
\theta_{\mathrm{aux}}=\Phi^\dagger G_\eta^{(0)},\qquad
P_\Phi=\Phi\Phi^\dagger,\qquad
D_\eta^{(0)}=G_\eta^{(0)}-\Phi\theta_{\mathrm{aux}}=(I_{Nn_{\mathrm y}}-P_\Phi)G_\eta^{(0)}.
\end{equation}
The single coefficient \(\theta_{\mathrm{aux}}\) is shared across all reference evaluations rather
than fitted separately at each one. The complementary range of the projector is the null space of
the transpose, which gives the orthogonality of the
construction~\cite[Lemma~3, Eq.~(11)]{gyorok2026obc}:
\begin{equation}\label{eq:originalorthogonality}
\boxed{\Phi^\top D_\eta^{(0)}=0.}
\end{equation}
The condition protects every baseline output direction reachable by a parameter vector, and it
does so as an aggregate over the reference set,
\begin{equation*}
(\Phi b)^\top D_\eta^{(0)}=0\ \ \text{for every }b\in\R^{n_{\theta_{\mathrm b}}},
\qquad\sum_{i=1}^N\phi_i^\top d_i^{(0)}=0,
\end{equation*}
so the summed products cancel while individual products need not vanish. It is therefore stronger
than orthogonality to a single baseline vector at one chosen parameter value.

The corrected predictor subtracts pointwise the same quantity that the stack subtracts on the
reference set,
\begin{equation}\label{eq:originalcorrected}
d_\eta^{(0)}(s,u)=g_{y,\eta}^{(0)}(s,u)-\phi(s,u)\theta_{\mathrm{aux}},\qquad
\hat y_k=\phi(s_k,u_k)\theta_b+d_\eta^{(0)}(s_k,u_k).
\end{equation}
After fitting, the coefficient is frozen at its estimated value, so that deployment evaluates
\begin{equation*}
\hat y_k=\phi(s_k,u_k)\hat\theta_b+g_{y,\hat\eta}^{(0)}(s_k,u_k)
-\phi(s_k,u_k)\hat\theta_{\mathrm{aux}}.
\end{equation*}
Evaluation is causal given the current regressor, and exact aggregate orthogonality holds on the
set from which the coefficient was constructed.

\section{The delay-state realisation}\label{sec:realisation}
The history \eqref{eq:history} is a state, and propagating it needs only fixed matrices whose
entries are zero or one. With shift matrices
\begin{equation*}
(S_{n_{\mathrm a}})_{ij}=\begin{cases}1&i=j+1,\\0&\text{otherwise},\end{cases}
\qquad S_{n_{\mathrm a}}\in\R^{n_{\mathrm a}\times n_{\mathrm a}},
\end{equation*}
with \(S_{n_{\mathrm b}}\) defined in the same way, and with first coordinate vectors of the
indicated lengths, set
\begin{equation}\label{eq:routing}
A=\diag(S_{n_{\mathrm a}}\otimes I_{n_{\mathrm y}},S_{n_{\mathrm b}}\otimes I_{n_{\mathrm u}}),
\qquad
B_x=\begin{bmatrix}e_1^{(n_{\mathrm a})}\otimes I_{n_{\mathrm y}}\\
0_{n_{\mathrm b}n_{\mathrm u}\times n_{\mathrm y}}\end{bmatrix},
\qquad
B_u=\begin{bmatrix}0_{n_{\mathrm a}n_{\mathrm y}\times n_{\mathrm u}}\\
e_1^{(n_{\mathrm b})}\otimes I_{n_{\mathrm u}}\end{bmatrix},
\end{equation}
so that the history update for an output sample \(v_k\) to be stored is
\begin{equation}\label{eq:storage}
s_{k+1}=As_k+B_uu_k+B_xv_k.
\end{equation}
The shift moves older samples down one slot. The two insertion matrices independently write the
current input and the new output. Their disjoint routing gives the three identities on which
everything below rests,
\begin{equation}\label{eq:routingidentities}
B_x^\top B_x=I_{n_{\mathrm y}},\qquad B_x^\top A=0,\qquad B_x^\top B_u=0,
\end{equation}
which hold because the newest-output rows of the shift and of the input insertion are zero, while
the output insertion copies each output coordinate exactly once. For scalar signals with two
output delays and one input delay the construction and its effect are
\begin{equation*}
s_k=\begin{bmatrix}y_{k-1}\\y_{k-2}\\u_{k-1}\end{bmatrix},\quad
A=\begin{bmatrix}0&0&0\\1&0&0\\0&0&0\end{bmatrix},\quad
B_x=\begin{bmatrix}1\\0\\0\end{bmatrix},\quad
B_u=\begin{bmatrix}0\\0\\1\end{bmatrix},\quad
s_{k+1}=\begin{bmatrix}v_k\\y_{k-1}\\u_k\end{bmatrix}.
\end{equation*}

The stored output sample admits two readings, which must be distinguished because they define
different predictors:
\begin{equation}\label{eq:feedback}
v_k=\begin{cases}\hat y_k&\text{recursive deterministic simulation},\\
y_k&\text{measured-history prediction}.\end{cases}
\end{equation}
In the recursive reading, history symbols denote values stored by the recursion itself once an
initial history has been supplied, and the baseline realisation is
\begin{equation}\label{eq:baselineSS}
\boxed{\begin{aligned}
s_{k+1}&=As_k+B_uu_k+B_x\phi(s_k,u_k)\theta_b,\\
\hat y_k&=\phi(s_k,u_k)\theta_b.
\end{aligned}}
\end{equation}
Selecting the newest slot returns the previously stored output rather than the current one,
\begin{equation*}
B_x^\top s_k=v_{k-1},\qquad B_x^\top s_{k+1}=v_k,
\end{equation*}
so the current output is supplied by the predictor function and never read out of the state. The
same routing carries the corrected predictor of \eqref{eq:originalcorrected}, provided the
corrected output appears both in the output equation and in the history feedback:
\begin{equation*}
s_{k+1}=As_k+B_uu_k+B_x\hat y_k,\qquad
\hat y_k=\phi(s_k,u_k)\theta_b+d_\eta^{(0)}(s_k,u_k).
\end{equation*}
Storing the predicted output already makes the source model recurrent under simulation, so the
learned states of Section~\ref{sec:memory} supply additional memory rather than the first memory.
The inputs may be prescribed or generated by a common causal controller.

In the measured reading the stored sample is the observation, which changes the feedback path but
not the routing,
\begin{equation*}
\hat y_k=\phi(s_k,u_k)\theta_b+d_\eta^{(0)}(s_k,u_k),\qquad
s_{k+1}=As_k+B_uu_k+B_xy_k.
\end{equation*}
For a declared generative process the same storage algebra applies to the data-generating signals
themselves,
\begin{equation*}
y_k=H(s_k,u_k)+e_k,\qquad
s_{k+1}=As_k+B_uu_k+B_xH(s_k,u_k)+B_xe_k,
\end{equation*}
which places the noise in the newest-output slot and nowhere else. This is bookkeeping for that
declared process and does not assert that any fitted predictor generates the measured data.

Realisations of this shape are standard for neural NARX models, and the routing matrices with
zero-one entries appear in Bonassi et al.~\cite[Sect.~II]{bonassi2022offsetfree} and in Seel et
al.~\cite[Sect.~III]{seel2021nnmpciss}, the latter also stating the feedback distinction of
\eqref{eq:feedback}. Their predictors index the output at the next sample while
\eqref{eq:baselineSS} keeps the current prediction index, and shifting the output index alone
would shift the last input index with it, so \eqref{eq:routing} is built for the timing used here
rather than obtained from theirs by reindexing. That the recursion and the realisation agree
exactly is established in Section~\ref{sec:projection}, once the model class is complete.

\section{Adding learned memory}\label{sec:memory}
The learning component of Section~\ref{sec:inherited} is a static map of the regressor, so its
correction at sample \(k\) is fixed by the stored history and the current input. We enlarge the
model class with an internal state carrying its own update, so that the correction may depend on
information the finite history does not hold:
\begin{equation}\label{eq:extendedIO}
\begin{aligned}
\hat y_k&=\phi(s_k,u_k)\theta_b+g_{y,\eta}(s_k,u_k,x_{a,k}),\\
x_{a,k+1}&=g_{a,\eta}(s_k,u_k,x_{a,k}).
\end{aligned}
\end{equation}
The learned parameters \(\eta\) are shared by the two maps, which act between
\begin{equation*}
g_{y,\eta}:\R^{n_s+n_{\mathrm u}+n_{x_a}}\to\R^{n_{\mathrm y}},\qquad
g_{a,\eta}:\R^{n_s+n_{\mathrm u}+n_{x_a}}\to\R^{n_{x_a}}.
\end{equation*}
Both read the old states and the current input rather than the same-step predicted output, so the
pair carries no algebraic loop. Memory is learned through the update in \eqref{eq:extendedIO}
instead of being stored as further signal delays, and the baseline predictor is untouched.

The learned state requires an initial value, either a fixed vector or the output of a causal
encoder \(\mathcal E_\eta\) evaluated on declared initial data:
\begin{equation*}
x_{a,0}=\bar x_a\qquad\text{or}\qquad x_{a,0}=\mathcal E_\eta(\text{declared initial data}).
\end{equation*}
Collecting the stored history and the learned state into the enlarged state
\begin{equation*}
\xi_k=\col(s_k,x_{a,k})\in\R^{n_s+n_{x_a}}
\end{equation*}
and substituting the output equation of \eqref{eq:extendedIO} into the routing \eqref{eq:storage}
separates the baseline and the learned transition contributions:
\begin{equation}\label{eq:extendedSS}
\boxed{\begin{aligned}
\xi_{k+1}&=F_b(\xi_k,u_k;\theta_b)+F_a(\xi_k,u_k;\eta),\\
F_b(\xi,u;\theta_b)&=\begin{bmatrix}As+B_uu+B_x\phi(s,u)\theta_b\\0_{n_{x_a}}\end{bmatrix},
\qquad
F_a(\xi,u;\eta)=\begin{bmatrix}B_xg_{y,\eta}(s,u,x_a)\\g_{a,\eta}(s,u,x_a)\end{bmatrix},\\
\hat y_k&=\phi(s_k,u_k)\theta_b+g_{y,\eta}(s_k,u_k,x_{a,k}).
\end{aligned}}
\end{equation}
The zero lower block is an assignment of the decomposition rather than a claim that the learned
coordinates stay at zero: a residual form of the learned update is represented inside
\(g_{a,\eta}\), and declaring a known latent baseline instead would move the split between
\(F_b\) and \(F_a\). Section~\ref{sec:projection} uses that zero block throughout.

\section{The extended projection}\label{sec:projection}
The orthogonality definition is applied directly to the two contributions to the enlarged state
equation. At the \(i\)-th stacked model evaluation, abbreviate
\(\phi_i=\phi(s_i,u_i)\),
\(g_{y,\eta,i}=g_{y,\eta}(s_i,u_i,x_{a,i})\), and
\(g_{a,\eta,i}=g_{a,\eta}(s_i,u_i,x_{a,i})\). Here \(x_{a,i}\) is the learned state already present
in the model evaluation; it is not a prescribed reference value or an additional item of training
data. The baseline and learned contributions are
\begin{equation*}
K_i=\begin{bmatrix}B_x\phi_i\\0_{n_{x_a}\times n_{\theta_{\mathrm b}}}\end{bmatrix},
\qquad
w_{\eta,i}=\begin{bmatrix}B_xg_{y,\eta,i}\\g_{a,\eta,i}\end{bmatrix},
\end{equation*}
where \(K_i\) is the partial derivative of the baseline transition with respect to the baseline
parameters at fixed arguments and not a total rollout derivative. Stacking these and factoring the
insertion out of the basis,
\begin{equation}\label{eq:augmentedstacks}
E=\begin{bmatrix}B_x\\0_{n_{x_a}\times n_{\mathrm y}}\end{bmatrix},\qquad
M_N=I_N\otimes E,\qquad
\Psi=\col_{i=1}^NK_i=M_N\Phi,\qquad
\mathcal G_\eta=\col_{i=1}^Nw_{\eta,i}.
\end{equation}
Define the output-correction stack at these same model evaluations by
\begin{equation*}
G_{y,\eta}=\col_{i=1}^N g_{y,\eta,i},\qquad
\theta_{\mathrm{aux}}=\Phi^\dagger G_{y,\eta}.
\end{equation*}
No separate set of learned-state reference values is introduced by these definitions.

\begin{proposition}[Zero-block reduction]\label{prop:zero}
In the stated Euclidean coordinates,
\begin{equation*}
E^\top E=I_{n_{\mathrm y}},\qquad M_N^\top M_N=I_{Nn_{\mathrm y}},\qquad
M_N^\top\mathcal G_\eta=G_{y,\eta},
\end{equation*}
and consequently
\begin{equation}\label{eq:zeroreduction}
\Psi^\top\mathcal G_\eta=\Phi^\top G_{y,\eta}.
\end{equation}
\end{proposition}

\begin{proof}
The first identity combines the output-insertion isometry of \eqref{eq:routingidentities} with the
zero lower block of \(E\), and the second follows from the first by the Kronecker product. At each
tuple the extraction reads
\begin{equation*}
E^\top w_{\eta,i}=B_x^\top B_xg_{y,\eta,i}+0^\top g_{a,\eta,i}=g_{y,\eta,i},
\end{equation*}
and stacking gives the third identity. Multiplying it by the baseline stack gives
\eqref{eq:zeroreduction}.
\end{proof}

\begin{lemma}[Isometric pseudoinverse]\label{lem:isometry}
With \(M_N\) as in \eqref{eq:augmentedstacks} and \(\Phi\) of any rank,
\begin{equation}\label{eq:isometricinverse}
(M_N\Phi)^\dagger=\Phi^\dagger M_N^\top.
\end{equation}
\end{lemma}

\begin{proof}
Let the rank \(\rho\) be positive and take a compact singular value decomposition, so that
\begin{equation*}
\Phi=U_\rho\Sigma_\rho V_\rho^\top,\qquad
(M_NU_\rho)^\top(M_NU_\rho)=I_\rho.
\end{equation*}
The insertion therefore preserves the nonzero singular values and supplies orthonormal left
singular vectors, whence
\begin{equation*}
(M_N\Phi)^\dagger=V_\rho\Sigma_\rho^{-1}U_\rho^\top M_N^\top=\Phi^\dagger M_N^\top.
\end{equation*}
At zero rank both pseudoinverses are zero and the identity holds trivially.
\end{proof}

Proposition~\ref{prop:zero} and Lemma~\ref{lem:isometry} together give the coefficient equivalence
at every rank:
\begin{equation}\label{eq:coefficientidentity}
\boxed{\Psi^\dagger\mathcal G_\eta=\Phi^\dagger G_{y,\eta}=\theta_{\mathrm{aux}}.}
\end{equation}
The enlarged projection therefore requires no additional condition on the learned-state update
in these coordinates. Its output/history block is projected against the same baseline parameter
span, while its additional-state rows are absent from that span. The learned state can still
influence the evaluated output correction and thereby the auxiliary coefficient.

The subtraction defined by that coefficient is applied to the output correction,
\begin{equation*}
d_\eta(s,u,x_a)=g_{y,\eta}(s,u,x_a)-\phi(s,u)\theta_{\mathrm{aux}},
\end{equation*}
and the corrected model and its realisation are
\begin{equation}\label{eq:finalIO}
\boxed{\begin{aligned}
\hat y_k&=\phi(s_k,u_k)\theta_b+d_\eta(s_k,u_k,x_{a,k}),\\
x_{a,k+1}&=g_{a,\eta}(s_k,u_k,x_{a,k}),
\end{aligned}}
\qquad
\boxed{\begin{aligned}
\xi_{k+1}&=F_b(\xi_k,u_k;\theta_b)+\widetilde F_a(\xi_k,u_k;\eta),\\
\widetilde F_a(\xi,u;\eta)&=\begin{bmatrix}B_xd_\eta(s,u,x_a)\\
g_{a,\eta}(s,u,x_a)\end{bmatrix}.
\end{aligned}}
\end{equation}
Realising this pair changes nothing that it predicts.

\begin{proposition}[Exact realisation]\label{prop:equivalence}
For matching initial histories, matching learned-state initialisation, matching parameters,
matching input rule and a matching auxiliary coefficient or causal coefficient schedule, the
extended input-output recursion \eqref{eq:finalIO} and its delay-state realisation generate
identical predictions and identical states at every finite sample at which the evaluations are
defined.
\end{proposition}

\begin{proof}
The stored histories and learned states coincide initially. Assume they coincide at a given
sample. Both forms then evaluate the same baseline, the same corrected output function and the
same learned update at the same arguments, so the predicted output and the next learned state
coincide, and \eqref{eq:storage} writes that output and the current input into precisely the
entries that \eqref{eq:history} requires at the next sample. Induction on the sample index
completes the claim. A common causal controller supplies equal inputs whenever its initialisation
and information agree.
\end{proof}

No rank, stability or observability hypothesis enters Proposition~\ref{prop:equivalence} beyond
the stated initialisation, timing, feedback and well-definedness, and the same induction covers
the model before subtraction. Under the measured reading of \eqref{eq:feedback} both forms store
observations instead,
\begin{equation*}
\hat y_k=\phi(s_k,u_k)\theta_b+d_\eta(s_k,u_k,x_{a,k}),\qquad
s_{k+1}=As_k+B_uu_k+B_xy_k,\qquad
x_{a,k+1}=g_{a,\eta}(s_k,u_k,x_{a,k}),
\end{equation*}
and for a declared generative process with a latent state the storage algebra of
Section~\ref{sec:realisation} extends unchanged,
\begin{equation*}
y_k=H(s_k,u_k,x_{a,k})+e_k,\qquad
\xi_{k+1}=\begin{bmatrix}As_k+B_uu_k+B_xH(s_k,u_k,x_{a,k})\\g_{a,\eta}(s_k,u_k,x_{a,k})\end{bmatrix}
+\begin{bmatrix}B_xe_k\\0_{n_{x_a}}\end{bmatrix},
\end{equation*}
which again confines the noise to the newest-output slot.

\begin{proposition}[Exact extended orthogonality]\label{prop:projection}
For the stacked evaluations used to construct the coefficient, the corrected learned transition
stack satisfies
\begin{equation}\label{eq:finalidentity}
\boxed{\begin{aligned}
\widetilde{\mathcal G}_\eta&=\mathcal G_\eta-\Psi\theta_{\mathrm{aux}}
=(I-\Psi\Psi^\dagger)\mathcal G_\eta,\\
\Psi^\top\widetilde{\mathcal G}_\eta
&=\Phi^\top(G_{y,\eta}-\Phi\Phi^\dagger G_{y,\eta})=0.
\end{aligned}}
\end{equation}
\end{proposition}

\begin{proof}
Equation \eqref{eq:coefficientidentity} identifies the state-space projector with the
output-coordinate one. Blockwise the subtraction and its orthogonality read
\begin{equation*}
\widetilde{\mathcal G}_\eta
=\col_i\begin{bmatrix}B_x(g_{y,\eta,i}-\phi_i\theta_{\mathrm{aux}})\\g_{a,\eta,i}\end{bmatrix},
\qquad\Phi^\top(I-\Phi\Phi^\dagger)=0,
\end{equation*}
and Proposition~\ref{prop:zero} completes the proof.
\end{proof}

Only the history rows receive a subtraction, so the learned update is numerically unchanged by the
projection at fixed parameters and arguments, while all dependence of the current output
correction on the learned state does participate. Its parameters can still move during joint
optimisation, through shared parameters and through the prediction loss.

A stronger statement holds for this particular routing. Writing the parameter-free part of the
baseline transition at the \(i\)-th stacked evaluation as
\begin{equation*}
b_i=\begin{bmatrix}As_i+B_uu_i\\0_{n_{x_a}}\end{bmatrix},\qquad
F_{b,i}=b_i+K_i\theta_b,
\end{equation*}
the identities \eqref{eq:routingidentities} make its inner product with each corrected learned
write vanish, so that the complete baseline transition is orthogonal to the complete corrected
learned transition:
\begin{equation}\label{eq:fullbaseline}
b_i^\top\widetilde F_{a,i}=0,\qquad
\sum_{i=1}^NF_{b,i}^\top\widetilde F_{a,i}=\theta_b^\top\Phi^\top D_{y,\eta}=0,
\qquad D_{y,\eta}=G_{y,\eta}-\Phi\theta_{\mathrm{aux}}.
\end{equation}
The primary projection \eqref{eq:finalidentity} protects the baseline parameter span, whereas
\eqref{eq:fullbaseline} additionally annihilates the parameter-free bookkeeping term and is
specific to the displayed routing.

\begin{remark}[Any rank]\label{rem:rank}
Nothing above assumes \eqref{eq:rank}. The Moore-Penrose convention selects the minimum-norm
least-squares coefficient; other minimisers differ by a null-space vector and agree on the
reference set while differing away from it,
\begin{equation*}
\theta'=\theta_{\mathrm{aux}}+v,\quad\Phi v=0,\qquad
\Phi\theta'=\Phi\theta_{\mathrm{aux}},\qquad
d'_\eta(s,u,x_a)-d_\eta(s,u,x_a)=-\phi(s,u)v.
\end{equation*}
One reference row and one unseen row already separate the two minimisers:
\begin{equation*}
\Phi=\begin{bmatrix}1&1\end{bmatrix},\qquad
v=\begin{bmatrix}1\\-1\end{bmatrix},\qquad
\phi_*=\begin{bmatrix}1&0\end{bmatrix},\qquad
\Phi v=0,\qquad \phi_*v=1.
\end{equation*}
\end{remark}

\begin{remark}[Differentiating through the coefficient]\label{rem:gradients}
The coefficient is a function of the learned parameters through the correction stack, so it is
recomputed whenever they change. For evaluation arguments fixed while differentiating the learned
parameters, the projection differentiates cleanly and the orthogonality is inherited by the
partial derivative:
\begin{equation}\label{eq:derivatives}
\partial_\eta\theta_{\mathrm{aux}}=\Phi^\dagger\partial_\eta G_{y,\eta},\qquad
\partial_\eta D_{y,\eta}=(I-P_\Phi)\partial_\eta G_{y,\eta},\qquad
\Phi^\top\partial_\eta D_{y,\eta}=0.
\end{equation}
The last identity is an aggregate derivative orthogonality. Writing the derivative evaluations
as \(J_i\), its meaning is
\begin{equation*}
\sum_{i=1}^N\phi_i^\top J_i=0.
\end{equation*}
The zero-covariance proof of Gy\"or\"ok et al.~\cite[Thm.~18, Eqs.~(39)--(40)]{gyorok2026obc}
displays pointwise block diagonality, whereas the stacked construction establishes only the
summed cross block. That identity makes the corresponding unweighted summed Gram matrix
block-diagonal; unequal output-channel noise weights require a separately justified weighted
cross-block identity. The covariance scope is stated in Section~\ref{sec:limits}. If an evaluation
argument itself depends on \(\eta\), its dependence belongs in a total derivative;
Equation~\eqref{eq:derivatives} does not silently treat such an argument as fixed.
\end{remark}

\begin{proposition}[Reduction to the inherited construction]\label{prop:reduction}
Remove the learned coordinates and the learned memory equation,
\begin{equation}\label{eq:reduction}
n_{x_a}=0,\qquad g_{y,\eta}(s,u)=g_{y,\eta}^{(0)}(s,u).
\end{equation}
Then the predictor, the stacks, the coefficient, the subtraction and the frozen deployment
expression of Section~\ref{sec:inherited} are recovered unchanged, and under \eqref{eq:rank} the
pseudoinverse coincides with the normal-matrix inverse of~\cite[Eq.~(8)]{gyorok2026obc}.
\end{proposition}

\begin{proof}
At \eqref{eq:reduction} the insertion of \eqref{eq:augmentedstacks} loses its zero block, so
\begin{equation*}
E=B_x,\qquad M_N=I_N\otimes B_x,\qquad \Psi=M_N\Phi,\qquad
\mathcal G_\eta=M_NG_\eta^{(0)},
\end{equation*}
and Proposition~\ref{prop:zero} with Lemma~\ref{lem:isometry} reduce
\eqref{eq:coefficientidentity} to the coefficient of \eqref{eq:originalprojection} and
\eqref{eq:finalidentity} to \eqref{eq:originalorthogonality}.
\end{proof}

This is a reduction of model functions and not a new statistical theorem.

\section{What is and is not established}\label{sec:limits}
The learned state reaches the current prediction through the correction, and both the next learned
state and the newly stored output reach later predictions,
\begin{equation*}
\begin{aligned}
x_{a,k}&\longrightarrow g_{y,\eta}(s_k,u_k,x_{a,k})\longrightarrow\hat y_k\longrightarrow s_{k+1},\\
x_{a,k}&\longrightarrow x_{a,k+1}.
\end{aligned}
\end{equation*}
This is compatible with \eqref{eq:finalidentity}, which concerns stacked state-function
evaluations and not propagated output responses. Orthogonality of accumulated output effects is a
different property, and it is not imposed here.

The projection identity does not transfer the source's parameter-recovery, consistency or
covariance theorems, and does not identify unique latent coordinates. It establishes nothing about
the empirical benefit of adding memory. In particular, the derivative identity establishes
aggregate unweighted orthogonality, not pointwise or general noise-weighted orthogonality.
Even pointwise unweighted orthogonality does not generally survive unequal channel weights;
the noise-weighted cross block in the source's covariance formula needs its own justification.
This qualifies the proof's reasoning without establishing that every conclusion of its theorem
is false. Stability of
the coupled history--learned-state predictor requires separate
analysis or an applicable stability certificate. A stable-by-design parametrisation is one
possible route, discussed by Gy\"or\"ok et al.~\cite[Remark~13]{gyorok2026obc}, who identify extension of
the cited parametrisations to augmentation as future work. Stability of an isolated learned-state
update does not automatically establish stability of the coupled predictor. The orthogonality of
the learned update block to zero in
\eqref{eq:finalidentity} is not a constraint on memory values or memory dynamics. The recursive
reading of \eqref{eq:feedback} is bookkeeping and implies no closed-loop stability property; the
stability theorems of the cited model predictive control work are not used here, and the routing
is constructed for the timing of \eqref{eq:baselineSS} rather than inherited from those sources.

The proofs use the direct-sum Euclidean geometry of \eqref{eq:augmentedstacks}. Block-diagonal
metrics compatible with that splitting retain the zero overlap on the learned block, but a
weighted history metric requires a matching weighted projection, cross blocks between history and
learned coordinates would let learned-memory writes enter the condition, and rescaling coordinates
requires transforming the metric to preserve the geometry. Those alternatives are not required by
the construction and are not analysed. Equation \eqref{eq:fullbaseline} is tied to that routing
more tightly still: an arbitrary physical baseline need not have its parameter-free transition
term orthogonal to the correction at all.

A rank deficiency in \eqref{eq:rank} on the stacked evaluations need not be a structural rank deficiency
of the function class, and the pseudoinverse does not itself resolve baseline parameter ambiguity,
as Remark~\ref{rem:rank} shows away from the reference set. A truncated singular value
decomposition protects the retained span exactly while dropping small nonzero singular values
gives only approximate orthogonality, and ridge inverses do not form exact orthogonal projectors.
Recomputing the coefficient numerically while stopping its gradient preserves the forward
orthogonal value but supplies a different derivative than \eqref{eq:derivatives}, whose identities
concern the construction evaluations and not free-run output sensitivity or covariance. The
zero-block reduction introduces no prescribed, frozen or separately generated learned-state
reference values. If histories, learned states or the basis depend on trainable quantities, their
dependence must be handled consistently when making derivative claims; the algebraic zero-block
identity itself remains a statement about the displayed state-equation blocks. Pseudoinverse
differentiation at rank changes needs separate care. Zero-history cases are also excluded from
Proposition~\ref{prop:equivalence}: an absent input history is handled by empty input blocks and
zero input insertion, whereas an absent output history requires separate treatment to retain the
insertion isometry.

Finally, nothing here establishes physical equivalence for a mechanical plant, encoder accuracy,
estimator consistency, covariance separation or closed-loop stability. What is established is the
additional-state extension of the model class and the exact reduction of its reference-set
projection in delay coordinates.

\appendix
\section{Symbolic verification}\label{sec:verification}
Selected identities were checked on a polynomial surrogate with scalar output, two output delays,
one input delay, two learned states and shared learned parameters. The 14 September 2026 run in
MATLAB R2025a Update 1 with Symbolic Math Toolbox 25.1 recorded eighteen passed checks,
reproduced in Table~\ref{tab:checks}. The verification records were inspected for this revision;
the computation was not rerun. These checks support the general proofs rather than verifying
every claim or every possible learned architecture.

\begin{table}[ht]
\centering\small
\begin{tabular}{llc}\toprule
\# & Check & Outcome\\\midrule
1 & Arbitrary one-step input-output and state-space transition agree & pass\\
2 & Output timing is the newest slot of the next history & pass\\
3 & Unprojected additional-state model & pass\\
4 & Shift rows orthogonal to the learned write & pass\\
5 & Input-output and state-space auxiliary coefficients identical & pass\\
6 & Input-output stacked orthogonality & pass\\
7 & State-space stacked orthogonality & pass\\
8 & Full baseline transition orthogonality & pass\\
9 & Differentiated projection & pass\\
10 & Rank-deficient coefficient equivalence & pass\\
11 & Rank-deficient orthogonality & pass\\
12.1 to 12.4 & Free simulation output, history and latent equality, four steps & pass\\
13 & Measured-history predictor equivalence & pass\\
14 & Stochastic history injection & pass\\
15 & Measured and predicted feedback are not interchangeable & pass\\
\bottomrule\end{tabular}
\caption{Symbolic checks and outcomes, eighteen checks in total.}\label{tab:checks}
\end{table}

Checks 1 and 13 compare symbolic one-step maps, while checks 12.1--12.4 unroll the recursion for
four steps. They support Proposition~\ref{prop:equivalence}, whose
induction supplies the general finite-sample argument. Physical equivalence for a mechanical plant
and the statistical theorems of the source were not tested.

% ================= STANDALONE WRAPPER (closing), delete with the opening ====
% CONVENTION (user-specified): BibTeX with a companion .bib file.
\bibliographystyle{plainnat}
\bibliography{obc-io-additional-state-chapter}
\end{document}
% ======================= END OF STANDALONE WRAPPER =========================
